\documentclass[12pt]{article}
\usepackage{amsmath, amssymb, amsthm}
\usepackage[english]{babel}
\usepackage{geometry}
\usepackage{mathrsfs}

\numberwithin{equation}{section}

% Other packages...
\newtheorem{theorem}{Theorem}[section]
\newtheorem{corollary}{Corollary}[theorem]
\newtheorem{lemma}[theorem]{Lemma}

\title{\textbf{Honors in Physics Research Report}}
\author{\textbf{Student}: Elisha Shmalo, \textbf{Mentor}: Dr. Jedediah H. Pixley}
\date{Date: \today}
\date{\textbf{Expected Graduation}: May 2026}
\date{\textbf{Reaserch Project Title}: Universality class of control transition in Heizenberg Spin Chain}

\geometry{left=1in, right=1in, top=1in, bottom=1in}
\renewcommand{\baselinestretch}{1.25}

\begin{document}

\maketitle

\textbf{I have emailed the final version of this report to my mentor:} We are studying a phase transition between overall 
chaotic and controlled dynamics in a controlled Heizenberg spin chain. We have coompleted the process of 
finding the critical exponents of the phase transition. Though we expected the transition to be of the Directed Percolation (DP)
universality class, the critical exponents we found differ from those of (DP). This was suprising as we had 
good reason to expect our transition to be DP. We recall now with seriousness our hamiltonian

\begin{equation}
    H = \sum_j (-\vec{J}(t) \circ \vec{S}_{j+1}) \cdot \vec{S}_j,
\end{equation}

in particular emphisizing that we made the $J_x, J_y$ coupling constants random in time (this was done to eliminate the 
presense of solitons that appeared in the spin chain). This time randomness is a none trivial pertubation in the 
sense of the Renormalization Group. It follows that it may be the reason our phase transition does not excibit the 
know critical values of DP. We are now seeking to dertermine 
if our phase transition belongs to a so called ``Time Random Directed Percolation" universality class. To do this we are 
preturbing (randomly in time) the control parameter of Stavskaya's probabilistic cellular automaton (STPCA). STPCA is known to 
be of the DP universality class. If the critical exponents of our random in time version of STPCA agree with the ones of 
our controlled spin chain, it will give us good reason to expect that the phase transition we are studing belongs to a 
``Time Random DP'' universality class. Finding a way to preturb STPCA to create time randomness similar to the kind in 
our spin chain (and such that the phase transition is preserved) has proven to be a challange, and one that is the subject of present reaserch.
In the end we plan to submit a manuscript of this work before the summer and we expect a publication.

\end{document}